\documentclass[11pt,a4paper]{ctexart}
\usepackage[margin=2.2cm]{geometry}
\usepackage{xcolor}
\usepackage{enumitem}
\usepackage{amssymb}
\usepackage[normalem]{ulem}
\usepackage{titlesec}
\usepackage{hyperref}
\definecolor{navy}{HTML}{1A365D}
\definecolor{ink}{HTML}{1F2933}
\definecolor{muted}{HTML}{66788A}
\definecolor{soft}{HTML}{EEF4FA}
\definecolor{gold}{HTML}{8A5A12}
\hypersetup{colorlinks=true,linkcolor=navy,urlcolor=navy}
\setlist{nosep}
\titleformat{\section}{\Large\bfseries\color{navy}}{}{0pt}{}
\titleformat{\subsection}{\large\bfseries\color{navy}}{}{0pt}{}
\pagestyle{plain}

\title{\textbf{七上名词与所有格专项}\\[4pt]
\large 人教版（新课纲版）七年级上册 · 词类 / 名词 / 所有格}
\author{初中英语资料站}
\date{2026-11-12}
\begin{document}
\maketitle

\section*{使用说明}
\begin{itemize}[leftmargin=1.4em]
\item 本讲义对应人教版（新课纲版）七年级上册语法要求：词类、名词与名词所有格。
\item 共 4 讲语法、6 级 48 题。
\item 在线版：\url{https://1901151458-prog.github.io/english-materials/g7a-nouns/}
\end{itemize}

\section{词类}
\begin{center}
\begin{tabular}{p{3.4cm}p{2.4cm}p{7.6cm}}
\hline
\textbf{词类} & \textbf{缩写} & \textbf{例词}\\
\hline
名词 & n. & son, tennis, family\\
冠词 & art. & a, an, the\\
代词 & pron. & we, her, that\\
形容词 & adj. & funny, different, excited\\
数词 & num. & nine, thirteen, first\\
动词 & v. & be, play, think\\
副词 & adv. & not, too, usually\\
介词 & prep. & in, on, at, of\\
连词 & conj. & and, but, because\\
感叹词 & interj. & well, hello, oh\\
\hline
\end{tabular}
\end{center}

\section{名词的种类与可数性}
\begin{itemize}[leftmargin=1.4em]
\item \textbf{专有名词}：特定的人、事件、机构、地点，如 Peter、January、Singapore，\textbf{第一个字母一般大写}。
\item \textbf{普通名词}：一般人或事物的名称，如 son、tennis、family。
\item \textbf{可数名词}有复数形式：eight rabbits、some trees。
\item \textbf{不可数名词}一般没有复数形式：milk、bread、yogurt。
\end{itemize}

\section{名词所有格：构成}
\begin{center}
\begin{tabular}{p{4.6cm}p{2.6cm}p{6.4cm}}
\hline
\textbf{类别} & \textbf{构成} & \textbf{例}\\
\hline
单数名词 & 加 's & Mike's basketball / Kate's cat / Alice's vlog\\
复数（不以 -s 结尾） & 加 's & Children's Day\\
复数（以 -s 结尾） & 加 ' & her grandparents' favourite grandchild\\
\hline
\end{tabular}
\end{center}
\textbf{口诀：}单数加 's；复数词尾有 s 只加 '；复数词尾无 s 仍加 's。

\section{名词所有格：读音}
\begin{center}
\begin{tabular}{p{5.4cm}p{2.6cm}p{5.2cm}}
\hline
\textbf{前一个音} & \textbf{'s 读} & \textbf{例}\\
\hline
清辅音 & /s/ & Mike's、Kate's\\
浊辅音、元音 & /z/ & your mum's、her dad's\\
/s/、/z/、/ʃ/、/tʃ/、/dʒ/ & /ɪz/ & Alice's、George's、the actress's\\
复数以 -s 结尾 & 读音不变 & grandparents'\\
\hline
\end{tabular}
\end{center}


\newpage
\section{分级练习（48 题）}
\noindent 按难度分级排列，先做基础，再做成分与扩展，最后做综合题。

\subsection*{第 1 级 · 词类与名词种类}
\noindent\textcolor{gray}{十大词类、专有名词与普通名词。}\par\smallskip
\begin{enumerate}[leftmargin=*,itemsep=4pt]
\item Peter、January、Singapore 属于哪类名词？
\begin{enumerate}[label=\Alph*.,leftmargin=2em,itemsep=1pt,topsep=2pt]
\item 专有名词
\item 普通名词
\item 不可数名词
\item 抽象名词
\end{enumerate}
\item son、tennis、family 属于哪类名词？
\begin{enumerate}[label=\Alph*.,leftmargin=2em,itemsep=1pt,topsep=2pt]
\item 专有名词
\item 普通名词
\item 数词
\item 代词
\end{enumerate}
\item 专有名词的第一个字母一般要 \_\_\_。
\begin{enumerate}[label=\Alph*.,leftmargin=2em,itemsep=1pt,topsep=2pt]
\item 大写
\item 小写
\item 斜体
\item 加粗
\end{enumerate}
\item 表示人或事物名称的词是 \_\_\_。
\begin{enumerate}[label=\Alph*.,leftmargin=2em,itemsep=1pt,topsep=2pt]
\item 名词
\item 动词
\item 形容词
\item 副词
\end{enumerate}
\item 表示动作或状态的词是 \_\_\_。
\begin{enumerate}[label=\Alph*.,leftmargin=2em,itemsep=1pt,topsep=2pt]
\item 名词
\item 动词
\item 形容词
\item 数词
\end{enumerate}
\item 用于修饰名词、表示人或事物特征的词是 \_\_\_。
\begin{enumerate}[label=\Alph*.,leftmargin=2em,itemsep=1pt,topsep=2pt]
\item 形容词
\item 副词
\item 介词
\item 连词
\end{enumerate}
\item 用于代替名词以及起名词作用的短语、分句或句子的词是 \_\_\_。
\begin{enumerate}[label=\Alph*.,leftmargin=2em,itemsep=1pt,topsep=2pt]
\item 代词
\item 数词
\item 冠词
\item 感叹词
\end{enumerate}
\item 表示数量或顺序的词是 \_\_\_。
\begin{enumerate}[label=\Alph*.,leftmargin=2em,itemsep=1pt,topsep=2pt]
\item 数词
\item 名词
\item 副词
\item 动词
\end{enumerate}
\end{enumerate}

\subsection*{第 2 级 · 可数与不可数}
\noindent\textcolor{gray}{可数名词有复数，不可数名词一般没有。}\par\smallskip
\begin{enumerate}[leftmargin=*,itemsep=4pt]
\item milk 属于哪类名词？
\begin{enumerate}[label=\Alph*.,leftmargin=2em,itemsep=1pt,topsep=2pt]
\item 可数名词
\item 不可数名词
\item 专有名词
\item 复数名词
\end{enumerate}
\item eight rabbits 中 rabbits 是 \_\_\_。
\begin{enumerate}[label=\Alph*.,leftmargin=2em,itemsep=1pt,topsep=2pt]
\item 可数名词复数
\item 不可数名词
\item 专有名词
\item 动词
\end{enumerate}
\item bread 属于哪类名词？
\begin{enumerate}[label=\Alph*.,leftmargin=2em,itemsep=1pt,topsep=2pt]
\item 可数名词
\item 不可数名词
\item 专有名词
\item 数词
\end{enumerate}
\item yogurt 属于哪类名词？
\begin{enumerate}[label=\Alph*.,leftmargin=2em,itemsep=1pt,topsep=2pt]
\item 可数名词
\item 不可数名词
\item 专有名词
\item 复数名词
\end{enumerate}
\item apple 的复数形式是 \_\_\_。
\begin{enumerate}[label=\Alph*.,leftmargin=2em,itemsep=1pt,topsep=2pt]
\item apples
\item apple’s
\item apples’
\item apple
\end{enumerate}
\item water 有复数形式吗？
\begin{enumerate}[label=\Alph*.,leftmargin=2em,itemsep=1pt,topsep=2pt]
\item 有
\item 没有
\item 只有 waters’
\item 只有 water’s
\end{enumerate}
\item some trees 中 trees 是 \_\_\_。
\begin{enumerate}[label=\Alph*.,leftmargin=2em,itemsep=1pt,topsep=2pt]
\item 可数名词复数
\item 不可数名词
\item 专有名词
\item 形容词
\end{enumerate}
\item book 的复数形式是 \_\_\_。
\begin{enumerate}[label=\Alph*.,leftmargin=2em,itemsep=1pt,topsep=2pt]
\item books
\item book’s
\item books’
\item book
\end{enumerate}
\end{enumerate}

\subsection*{第 3 级 · 名词所有格 · 形式}
\noindent\textcolor{gray}{单数加 ’s；复数词尾有 s 只加 ’。}\par\smallskip
\begin{enumerate}[leftmargin=*,itemsep=4pt]
\item Mike 的篮球：\_\_\_。
\begin{enumerate}[label=\Alph*.,leftmargin=2em,itemsep=1pt,topsep=2pt]
\item Mike’s basketball
\item Mikes basketball
\item Mike basketball
\item Mike’s basketball’s
\end{enumerate}
\item Children’s Day 中 ’s 加在 \_\_\_ 上。
\begin{enumerate}[label=\Alph*.,leftmargin=2em,itemsep=1pt,topsep=2pt]
\item 复数名词（不以 -s 结尾）
\item 单数名词
\item 复数名词（以 -s 结尾）
\item 动词
\end{enumerate}
\item 她祖父母的（所有格）：\_\_\_。
\begin{enumerate}[label=\Alph*.,leftmargin=2em,itemsep=1pt,topsep=2pt]
\item grandparents’
\item grandparents’s
\item grandparent’s
\item grandparents
\end{enumerate}
\item Kate 的猫：\_\_\_。
\begin{enumerate}[label=\Alph*.,leftmargin=2em,itemsep=1pt,topsep=2pt]
\item Kate’s cat
\item Kates cat
\item Kate cat
\item Kates’ cat
\end{enumerate}
\item Alice 的 vlog：\_\_\_。
\begin{enumerate}[label=\Alph*.,leftmargin=2em,itemsep=1pt,topsep=2pt]
\item Alice’s vlog
\item Alices vlog
\item Alice vlog
\item Alices’ vlog
\end{enumerate}
\item 复数名词以 -s 结尾，所有格加 \_\_\_。
\begin{enumerate}[label=\Alph*.,leftmargin=2em,itemsep=1pt,topsep=2pt]
\item ’
\item ’s
\item s’
\item s
\end{enumerate}
\item 单数名词所有格加 \_\_\_。
\begin{enumerate}[label=\Alph*.,leftmargin=2em,itemsep=1pt,topsep=2pt]
\item ’s
\item ’
\item s
\item 不变化
\end{enumerate}
\item 那位女演员的名字：\_\_\_。
\begin{enumerate}[label=\Alph*.,leftmargin=2em,itemsep=1pt,topsep=2pt]
\item the actress’s name
\item the actress’ name
\item the actresses name
\item the actress name
\end{enumerate}
\end{enumerate}

\subsection*{第 4 级 · 名词所有格 · 读音}
\noindent\textcolor{gray}{/s/、/z/、/ɪz/ 的规则。}\par\smallskip
\begin{enumerate}[leftmargin=*,itemsep=4pt]
\item Mike’s 中 ’s 读 \_\_\_。
\begin{enumerate}[label=\Alph*.,leftmargin=2em,itemsep=1pt,topsep=2pt]
\item /s/
\item /z/
\item /ɪz/
\item /t/
\end{enumerate}
\item your mum’s 中 ’s 读 \_\_\_。
\begin{enumerate}[label=\Alph*.,leftmargin=2em,itemsep=1pt,topsep=2pt]
\item /s/
\item /z/
\item /ɪz/
\item /d/
\end{enumerate}
\item Kate’s 中 ’s 读 \_\_\_。
\begin{enumerate}[label=\Alph*.,leftmargin=2em,itemsep=1pt,topsep=2pt]
\item /s/
\item /z/
\item /ɪz/
\item /eɪ/
\end{enumerate}
\item Alice’s 中 ’s 读 \_\_\_。
\begin{enumerate}[label=\Alph*.,leftmargin=2em,itemsep=1pt,topsep=2pt]
\item /s/
\item /z/
\item /ɪz/
\item /iː/
\end{enumerate}
\item George’s 中 ’s 读 \_\_\_。
\begin{enumerate}[label=\Alph*.,leftmargin=2em,itemsep=1pt,topsep=2pt]
\item /s/
\item /z/
\item /ɪz/
\item /dʒ/
\end{enumerate}
\item 复数名词以 -s 结尾时，所有格读音 \_\_\_。
\begin{enumerate}[label=\Alph*.,leftmargin=2em,itemsep=1pt,topsep=2pt]
\item 不变
\item 加 /z/
\item 加 /ɪz/
\item 加 /s/
\end{enumerate}
\item her dad’s 中 ’s 读 \_\_\_。
\begin{enumerate}[label=\Alph*.,leftmargin=2em,itemsep=1pt,topsep=2pt]
\item /s/
\item /z/
\item /ɪz/
\item /d/
\end{enumerate}
\item 在元音后，’s 通常读 \_\_\_。
\begin{enumerate}[label=\Alph*.,leftmargin=2em,itemsep=1pt,topsep=2pt]
\item /s/
\item /z/
\item /ɪz/
\item /eɪ/
\end{enumerate}
\end{enumerate}

\subsection*{第 5 级 · 词类辨析}
\noindent\textcolor{gray}{按例词判断属于哪一词类。}\par\smallskip
\begin{enumerate}[leftmargin=*,itemsep=4pt]
\item funny、different、excited 属于哪一词类？
\begin{enumerate}[label=\Alph*.,leftmargin=2em,itemsep=1pt,topsep=2pt]
\item 形容词
\item 副词
\item 名词
\item 动词
\end{enumerate}
\item in、on、at、of 属于哪一词类？
\begin{enumerate}[label=\Alph*.,leftmargin=2em,itemsep=1pt,topsep=2pt]
\item 介词
\item 连词
\item 冠词
\item 副词
\end{enumerate}
\item and、but、because、when 属于哪一词类？
\begin{enumerate}[label=\Alph*.,leftmargin=2em,itemsep=1pt,topsep=2pt]
\item 连词
\item 介词
\item 副词
\item 代词
\end{enumerate}
\item we、her、that、what 属于哪一词类？
\begin{enumerate}[label=\Alph*.,leftmargin=2em,itemsep=1pt,topsep=2pt]
\item 代词
\item 名词
\item 形容词
\item 数词
\end{enumerate}
\item nine、thirteen、first 属于哪一词类？
\begin{enumerate}[label=\Alph*.,leftmargin=2em,itemsep=1pt,topsep=2pt]
\item 数词
\item 代词
\item 副词
\item 冠词
\end{enumerate}
\item not、too、there、usually 属于哪一词类？
\begin{enumerate}[label=\Alph*.,leftmargin=2em,itemsep=1pt,topsep=2pt]
\item 副词
\item 介词
\item 连词
\item 形容词
\end{enumerate}
\item a、an、the 属于哪一词类？
\begin{enumerate}[label=\Alph*.,leftmargin=2em,itemsep=1pt,topsep=2pt]
\item 冠词
\item 代词
\item 数词
\item 介词
\end{enumerate}
\item well、hello、oh、hey 属于哪一词类？
\begin{enumerate}[label=\Alph*.,leftmargin=2em,itemsep=1pt,topsep=2pt]
\item 感叹词
\item 副词
\item 连词
\item 名词
\end{enumerate}
\end{enumerate}

\subsection*{第 6 级 · 综合提升}
\noindent\textcolor{gray}{名词、所有格与词类混合辨析。}\par\smallskip
\begin{enumerate}[leftmargin=*,itemsep=4pt]
\item 名词按可数性分为 \_\_\_。
\begin{enumerate}[label=\Alph*.,leftmargin=2em,itemsep=1pt,topsep=2pt]
\item 可数名词和不可数名词
\item 专有名词和普通名词
\item 单数和复数
\item 具体和抽象
\end{enumerate}
\item Singapore 这类专有名词首字母 \_\_\_。
\begin{enumerate}[label=\Alph*.,leftmargin=2em,itemsep=1pt,topsep=2pt]
\item 必须大写
\item 必须小写
\item 可大写可小写
\item 加引号
\end{enumerate}
\item the students’ books 表示 \_\_\_。
\begin{enumerate}[label=\Alph*.,leftmargin=2em,itemsep=1pt,topsep=2pt]
\item 学生们的书
\item 一个学生的书
\item 书的名称
\item 学生和书
\end{enumerate}
\item 单数名词所有格 \_\_\_。
\begin{enumerate}[label=\Alph*.,leftmargin=2em,itemsep=1pt,topsep=2pt]
\item 加 ’s
\item 只加 ’
\item 加 s
\item 不变
\end{enumerate}
\item 复数名词不以 -s 结尾时，所有格 \_\_\_。
\begin{enumerate}[label=\Alph*.,leftmargin=2em,itemsep=1pt,topsep=2pt]
\item 加 ’s
\item 只加 ’
\item 加 s’
\item 不变
\end{enumerate}
\item milk 是不可数名词，\_\_\_ 复数形式。
\begin{enumerate}[label=\Alph*.,leftmargin=2em,itemsep=1pt,topsep=2pt]
\item 没有
\item 有
\item 有两种
\item 只有 milk’s
\end{enumerate}
\item 表示顺序先后的数词叫 \_\_\_。
\begin{enumerate}[label=\Alph*.,leftmargin=2em,itemsep=1pt,topsep=2pt]
\item 序数词
\item 基数词
\item 数词
\item 名词
\end{enumerate}
\item 用于连接单词、短语或句子的词类是 \_\_\_。
\begin{enumerate}[label=\Alph*.,leftmargin=2em,itemsep=1pt,topsep=2pt]
\item 连词
\item 介词
\item 副词
\item 冠词
\end{enumerate}
\end{enumerate}

\newpage
\section*{答案与解析}
\small
\begin{enumerate}[leftmargin=1.6em,itemsep=2pt]
\item \textbf{N1-1 A}：特定的人、时间、地点的专有名称是专有名词；教材例句。
\item \textbf{N1-2 B}：表示一般人或事物的名称，是普通名词。
\item \textbf{N1-3 A}：教材明确：专有名词的第一个字母一般要大写。
\item \textbf{N1-4 A}：名词（Noun）表示人或事物的名称；表1。
\item \textbf{N1-5 B}：动词（Verb）表示动作或状态；表1。
\item \textbf{N1-6 A}：形容词（Adjective）用于修饰名词；表1。
\item \textbf{N1-7 A}：代词（Pronoun）；表1。
\item \textbf{N1-8 A}：数词（Numeral）；表1。
\item \textbf{N2-1 B}：milk 一般没有复数形式，是不可数名词；教材例句。
\item \textbf{N2-2 A}：rabbit 有复数形式，是可数名词；教材例句。
\item \textbf{N2-3 B}：bread 一般没有复数形式，是不可数名词；教材例句。
\item \textbf{N2-4 B}：yogurt 是不可数名词；教材例句。
\item \textbf{N2-5 A}：可数名词一般加 -s 构成复数。
\item \textbf{N2-6 B}：不可数名词一般没有复数形式。
\item \textbf{N2-7 A}：tree 有复数形式，是可数名词。
\item \textbf{N2-8 A}：可数名词一般加 -s。
\item \textbf{N3-1 A}：单数名词所有格加 ’s；教材表2。
\item \textbf{N3-2 A}：children 是复数但不以 -s 结尾，加 ’s；教材表2。
\item \textbf{N3-3 A}：复数名词以 -s 结尾，只加 ’；教材表2。
\item \textbf{N3-4 A}：单数名词加 ’s；教材表2。
\item \textbf{N3-5 A}：单数名词加 ’s；教材表2。
\item \textbf{N3-6 A}：只加一个撇号；教材表2。
\item \textbf{N3-7 A}：单数名词加 ’s；教材表2。
\item \textbf{N3-8 A}：单数名词加 ’s；教材表2 例句 the actress’s name。
\item \textbf{N4-1 A}：在清辅音 /k/ 后读 /s/；教材表2。
\item \textbf{N4-2 B}：在浊辅音 /m/ 后读 /z/；教材表2。
\item \textbf{N4-3 A}：在清辅音 /t/ 后读 /s/；教材表2。
\item \textbf{N4-4 C}：在 /s/ 音后读 /ɪz/；教材表2。
\item \textbf{N4-5 C}：在 /dʒ/ 后读 /ɪz/；教材表2。
\item \textbf{N4-6 A}：教材表2：读音不变。
\item \textbf{N4-7 B}：在浊辅音 /d/ 后读 /z/；教材表2。
\item \textbf{N4-8 B}：教材表2：在浊辅音和元音后读 /z/。
\item \textbf{N5-1 A}：形容词（adj.）；表1。
\item \textbf{N5-2 A}：介词（prep.）；表1。
\item \textbf{N5-3 A}：连词（conj.）；表1。
\item \textbf{N5-4 A}：代词（pron.）；表1。
\item \textbf{N5-5 A}：数词（num.）；表1。
\item \textbf{N5-6 A}：副词（adv.）；表1。
\item \textbf{N5-7 A}：冠词（art.）；表1。
\item \textbf{N5-8 A}：感叹词（interj.）；表1。
\item \textbf{N6-1 A}：教材：名词按其可数性可分为可数名词和不可数名词。
\item \textbf{N6-2 A}：专有名词的第一个字母一般要大写。
\item \textbf{N6-3 A}：复数以 -s 结尾加 ’，表示「……们的」。
\item \textbf{N6-4 A}：教材表2。
\item \textbf{N6-5 A}：如 Children’s Day；教材表2。
\item \textbf{N6-6 A}：不可数名词一般没有复数形式。
\item \textbf{N6-7 A}：教材：数词包括基数词和序数词，序数词表示顺序的先后。
\item \textbf{N6-8 A}：连词（Conjunction）；表1。
\end{enumerate}

\end{document}
